\documentclass[12pt]{article}
\usepackage[utf8]{inputenc}
\usepackage[T1]{fontenc}
\usepackage{hyperref}
\title{Dori'o Language Resources for K-12 Curriculum Development}
\author{Compiled by Claude Sonnet 5.5 (claude-sonnet-5-5) for Andr\'as Kornai}
\date{}
\begin{document}
\maketitle

\section{Introduction}
Dori'o (ISO 639-3 \texttt{dor}, also called Kwarekwareo) is an Austronesian language of South Malaita, Solomon Islands, with about 2,400 speakers cited from 1999. The search found no dedicated dictionary, grammar or school materials.

\section{Primary Educational Resources}
No primers, textbooks or curricula in Dori'o were found.

\section{Linguistic and Reference Materials}
\begin{itemize}
\item Dori'o language, Wikipedia overview: \url{https://en.wikipedia.org/wiki/Dori%CA%BCo_language}
\item Numerals from a 1988 elicitation by Darrel Tryon (MPI): \url{https://lingweb.eva.mpg.de/channumerals/Dori'o.htm}
\item Verbix verb list with glosses: \url{https://docs.verbix.com/Languages/Dorio}
\item Glosbe entry: \url{https://glosbe.com/nmx/dor}
\item Ivens, \emph{Dictionary and Grammar of the Language of Sa'a and Ulawa} (1918). This covers a neighbouring, related language, not Dori'o, and is listed only as comparative background: \url{https://bookserver.archive.org/acquire/dictionarygramma00ivenuoft.pdf}
\end{itemize}

\section{Religious and Community Resources}
None specific to Dori'o were found.

\section{Digital and Audio Resources}
No audio, video or software resources were found.

\section{Curriculum Development Context}
The language is reported as stable and not known to be taught in schools. Existing material is limited to short word lists and numerals, so a curriculum would depend on new documentation with the community.
\end{document}